\begin{resumo}
	\par A autenticação biométrica baseada em fala é comprometida quando o locutor apresenta disfonia severa, condição que altera substancialmente as características acústicas da voz. Este trabalho propõe complementar os sinais de voz fonada com sinais de eletroencefalograma (EEG) capturados durante a fala imaginada, fenômeno em que o indivíduo articula mentalmente sem emitir som, ativando regiões cerebrais análogas às da fala fonada. Duas estratégias de extração de características são comparadas: a manual, baseada na Transformada \textit{Wavelet-Packet} de Tempo Discreto (DTWPT) combinada com a Engenharia Paraconsistente de Características; e a automatizada, por meio de \textit{autoencoders} profundos. As características extraídas são classificadas por Redes Neurais de Pulso (RNP) em uma arquitetura residual, motivadas pela eficiência energética e pela capacidade nativa de processar sequências temporais. A abordagem é avaliada sobre a base pública de EEG para fala imaginada \cite{10.1117/12.2255697}, sob protocolo de verificação --- em que os locutores de teste não aparecem no treino --- tendo como métricas primárias a Taxa de Erro Igual (EER) e a área sob a curva ROC (AUC); o Erro Médio Quadrático é usado apenas na reconstrução dos \textit{autoencoders}. A seleção do extrator de características é feita pela engenharia paraconsistente, sem treinar classificador, e apenas a combinação vencedora alimenta a etapa de autenticação. Os resultados obtidos mostram que a extração manual supera as variantes baseadas em \textit{autoencoder} segundo o critério paraconsistente, e que o desempenho de autenticação permanece próximo do nível de acaso (melhor EER de $0{,}4534$), indicando que a fusão de voz e EEG, nas condições avaliadas, ainda não sustenta uma autenticação prática.
\end{resumo}

\begin{resumo}[Abstract]
	\begin{otherlanguage*}{english}
		\par Speech-based biometric authentication degrades severely when the speaker presents pronounced dysphonia, as this condition substantially alters the acoustic characteristics of the voice. This work proposes to complement phonated speech signals with electroencephalogram (EEG) signals captured during imagined speech — a phenomenon in which an individual mentally articulates without producing sound, activating brain regions analogous to those involved in phonated speech. Two feature-extraction strategies are compared: a manual approach based on the Discrete Time Wavelet Packet Transform (DTWPT) combined with Paraconsistent Feature Engineering; and an automated approach using deep autoencoders. The extracted features are classified by Spiking Neural Networks (SNNs) in a residual architecture, motivated by their energy efficiency and inherent ability to process temporal sequences. The approach is evaluated on the public EEG imagined-speech database \cite{10.1117/12.2255697} under a verification protocol --- in which test speakers do not appear in training --- with Equal Error Rate (EER) and area under the ROC curve (AUC) as primary metrics; Mean Squared Error is used only for autoencoder reconstruction. Feature-extractor selection is performed by Paraconsistent Feature Engineering, without training any classifier, and only the winning combination feeds the authentication stage. The obtained results show that handcrafted extraction outperforms the autoencoder variants under the paraconsistent criterion, and that authentication performance remains close to chance level (best EER of $0.4534$), indicating that the fusion of voice and EEG, under the evaluated conditions, does not yet support practical authentication.
	\end{otherlanguage*}
\end{resumo}
